\documentclass[11pt,a4paper]{article}
\usepackage[utf8]{inputenc}
\usepackage[T1]{fontenc}
\usepackage{amsmath,amssymb,amsthm}
\usepackage{geometry}
\geometry{margin=2.4cm}

\newtheorem{definition}{Definition}
\newtheorem{theorem}{Theorem}
\newtheorem{remark}{Remark}

\title{Logarithmic Scales ($\ln 2$, $\ln \phi$, and $\ln 10$) Across Thermodynamics, Information, and Geometry: An OVRP Observational and Formal Audit}
\author{Mariano Panzano Caball\'e \\ \small Oasis Sovereign Architecture Research, Barcelona}
\date{November 2026}

\begin{document}
\maketitle

\begin{abstract}
Within the \emph{Oasis Verifiable Research Pipeline} (OVRP, \texttt{DOI: 10.5281/zenodo.24443531}), we examine the roles of $\ln 2 \approx 0.693147$, the golden-ratio logarithm $\ln \phi \approx 0.481212$ ($\phi = \frac{1+\sqrt{5}}{2} \approx 1.618034$), and $\ln 10 \approx 2.302585$ across thermodynamics, information theory, and information geometry. Rather than claiming a universal law or closing open foundational problems, we document $\ln 10$ as an observed attractor/reference scale in our probes alongside $\ln 2$ and $\ln \phi$, contextualized against recent literature on information thermodynamics and geometric/entropic structures. A self-contained \texttt{Lean~4.36.0} specification (\texttt{Ln10Ln2\_TGI\_Spec.lean}, \texttt{0 sorry}, \texttt{0} custom axioms) and a deterministic \texttt{C99} probe accompany the manuscript under \texttt{SHA-256} and Bitcoin OpenTimestamps provenance.
\end{abstract}

\section{Thermodynamics, Information, Geometry, and Observed Scales}
In information thermodynamics and quantum information, $\ln 2$ marks the binary/two-state logarithmic scale, whereas $\ln 10$ provides a ten-base logarithmic scale ($\log_{10} x = \ln x / \ln 10$) observed as a recurring attractor/reference value across our OVRP probes. Incorporating $\ln \phi \approx 0.481212$ allows a clear comparative table across the three logarithmic scales ($1 - \ln \phi / \ln 2 \approx 30.58\%$ relative gap vs.\ $\ln 2$), keeping physical thermodynamic bounds distinct from structural/combinatorial comparisons.

\begin{table}[h!]
\centering
\caption{Comparative Logarithmic Scales ($\ln 2$, $\ln \phi$, and $\ln 10$).}
\begin{tabular}{cccc}
\hline
Constant / Scale & Value (nats) & Ratio to $\ln 2$ & Role in Comparative Audit \\
\hline
$\ln 2$ & $0.693147$ & $1.000000$ & Binary / single-qubit entropy reference \\
$\ln \phi$ & $0.481212$ & $0.694243$ & Golden-ratio log ($\phi \approx 1.618034$; $30.58\%$ below $\ln 2$) \\
$\ln 10$ & $2.302585$ & $3.321928$ & Decimal scale \& observed probe attractor/reference \\
\hline
\end{tabular}
\end{table}

\section{Conclusions and Open Boundaries}
\begin{itemize}
  \item \textbf{Scale Comparison ($\ln 2$, $\ln \phi$, $\ln 10$):} Putting $\ln 2$, $\ln \phi$, and $\ln 10$ side by side clarifies their exact numerical ratios without conflating base conversions or structural state-space ratios with a universal physical law.
  \item \textbf{Observed Attractor vs.\ Open Theory:} Documenting $\ln 10$ as an observed attractor/reference in our probes anchors the narrative transparently while leaving open questions in non-equilibrium quantum thermodynamics, information geometry, gauge confinement, and analytic number theory explicit.
\end{itemize}

\begin{remark}[Explicit Open Boundaries]
Observing $\ln 10$ as an attractor in specific probes and models does not establish a universal physical law or resolve open problems in quantum gravity, gauge confinement, or analytic number theory.
\end{remark}

\end{document}
